\section{Tool Overview and Semantic Model}

A test consists of an annotated C function, its logical dependencies, a typed
partial binding $\theta$, and an execution/checking policy. The binding may fix
scalar arguments, logical data, lengths, or constructor terms. Pointer addresses
can remain symbolic. The spatial precondition describes memory through the
ordinary QCP assertions; the controller does not construct a native heap.

Let $R$ and $E$ be the entry and return assertions. An initial state includes
the C environment and heap, and an interpretation of unbound logical variables.
The instance domain is
\[
  D_\theta = \{\sigma \mid \sigma \models R[\theta]\}.
\]
The entry check asks whether $D_\theta$ is nonempty. Execution checking is
universal over this domain. A witness used to establish nonemptiness does not
silently replace the instance by that witness. For example, fixing a list
length to two need not fix node addresses or every payload value.

The implementation records entry satisfiability, execution completion, safety,
return obligations and checks at intermediate boundaries separately. A completed
execution test passes only if the entry is satisfiable and all selected
obligations close. An unsatisfiable entry fails the entry check. A remaining VC,
execution budget exhaustion, or unsupported call leaves an inconclusive result.
A failed attempt to prove a VC establishes neither its negation nor a reachable
counterexample. Proof-certified results additionally require checked evidence
covering the relevant obligations and an explicit account of trusted assumptions.

Comparative evaluations must freeze the input domain independently of the
candidate under evaluation. Changing a candidate's Require can change
$D_\theta$ even when $\theta$ is unchanged. Accordingly, the experimental
protocol separates contract-driven interactive tests from matched, independently
specified input domains and reports entry rejection separately from functional
violations.

\section{Approach}

\subsection{Typed instantiation and entry satisfiability}

The analyzer identifies C parameters, typed \texttt{With} variables and imported
logical dependencies. Bindings denote a capture-avoiding substitution on logical
variable identities. Each call and quantified definition is freshened before
substitution. Logical inputs are immutable entry values even when the program
subsequently changes its parameters or reachable memory.

The instantiated Require is simplified before checking satisfiability.
The checker detects binding conflicts, impossible value ranges and overlapping
separately owned memory. A positive SMT result is sufficient only for a
faithfully encoded fragment; an abstraction containing necessary constraints
can establish inconsistency but cannot alone establish spatial satisfiability.
Unresolved spatial and logical requirements produce an existential Rocq goal
with the relevant definitions and layout information. Arbitrary spatial
predicates are never accepted merely because disjoint addresses seem available.

\subsection{One evolving logical state}

The execution configuration contains the current control location, QCP assertion
$P$, call stack, accumulated obligations and remaining budget. QCP evaluates
program expressions, reads and writes memory, and produces successor assertions.
Simplification runs on these assertions throughout execution. A determined
condition selects its feasible successor; a symbolic condition retains every
successor whose infeasibility has not been established. No external numerical
address or native execution is needed to choose an owned symbolic cell.

\subsection{Joint heap and value normalization}
\label{sec:normalization-algorithm}
\label{sec:normalization}
The continuation can refer to logical values outside the current heap formula.
Eliminating $r$ from $\exists r.\ r=2\land\mathit{cell}(p,1)$ while leaving a
return term that refers to $r$ loses its value constraint. The normalization
interface therefore preserves the assertion $P$ together with every live term
$\bar t$, including results and saved call values. Define
\[
\obs(P,\bar t)(s,\bar v)=\exists\eta.\ P(s,\eta)\land
                  \bar v=\operatorname{eval}(\bar t)_{s,\eta}.
\]
Each admitted rewrite preserves this relation. Its composition consequently
preserves the observations used by entry and boundary checks.

\begin{lstlisting}[language={},caption={Normalization of the state and its live values.},label={lst:normalize}]
Normalize(Gamma, P, liveTerms, budget):
  while budget > 0 and an admitted reduction exists:
    resolve its definition and interpretation
    establish its guard; freshen typed local witnesses
    replace the occurrence by the complete equivalent body
    update dependent liveTerms consistently
    simplify exposed computations
    budget := budget - 1
  return (P, liveTerms)
\end{lstlisting}

\paragraph{Admitting reductions.}
Closed computations and known constructors expose reducible expressions.
Definition unfolding retains every pure conjunct, spatial resource, and typed
witness. For a guarded clause, activation requires
$g(\bar a)\Rightarrow(H(\bar a)\equiv\exists\bar w.\,B(\bar a,\bar w))$
in the selected interpretation. The executor establishes $g$ from the current
state. Unresolved guards retain the predicate application. Budgeted rewriting
terminates while preserving the meaning of residual terms.

For optional external rules, the controller checks the equation and its
association with the program's qualified definitions, constructors, spatial
notation, and layout. This prevents a same-named imported definition from
changing the meaning of an accepted equation. Imported predicates otherwise
use QCP's existing strategies. A checked pure consequence can be exposed by
$H\equiv H\land G$, retaining the original spatial predicate and its witnesses.
This path supports local proof progress without requiring a full layout
translation.

\paragraph{Preserving witnesses and symbolic addresses.}
Existing live variables remain protected during expansion. An unused witness
of nominal type $T$ requires inhabitance:
$(\exists w:T.\,P)\equiv(\mathit{inhabited}(T)\land P)$ when $w$ is absent
from $P$. Integer and list types supply canonical witnesses; an arbitrary
nominal type retains its existence obligation. Shared-base offset reductions
and closed list slices simplify operations with symbolic pointers. Their
ordinary C bounds and arithmetic safety obligations remain in the VC interface.


The normalizer evaluates available closed expressions and unfolds definitions
when their arguments and guards permit it. The current implementation expands
frontend predicate definitions and supported executable logical operations,
and uses supplied entailment strategies. Staging an imported Rocq definition
for proof does not itself make that definition executable by the normalizer.
Other terms remain logical terms; an opaque heap predicate can therefore
prevent a required memory read even when its definition is available to ITP.
External definition expansion is optional and disabled by default. Ordinary
runs retain imported predicates and use QCP's supplied strategies. Availability
of a definition alone is not evidence of solver benefit; the controlled
comparison below separates predicate-preserving pure consequences from full
spatial expansion and counts complete test outcomes as well as residual VCs.
The candidate executor additionally accepts a dedicated scope of checked
definition equations. Its controller freezes the module identity, notation
environment and generated rules, then compiles the equations, audits assumptions
and replays them in the Rocq kernel. Activation pins a dedicated file stem and
the certified left-rule IDs; unrelated rules in the same scope are excluded,
and duplicate file/ID identities are rejected. The executor uses QCP's matcher to apply
only left-side unfolding actions at state updates and precondition checking;
goal-side model proposals remain proof obligations. The explicit controller
checks both the equations and their association with a selected program instance.
A parse-only preflight exports the actual layout and Rocq import context.
The association certificate captures the program's interpretation of each
referenced predicate and constructor before importing the equation context,
then establishes identity with the corresponding frozen equation bindings.
The generated equation context imports the selected definition's declared
dependencies before the definition module. The equation's left-hand side,
recursive occurrences and model-choice statements use the selected predicate's
fully qualified name. Writing $f=M.p$ for that definition, $b_p$ for its frozen
equation binding and $c_p$ for the program binding captured before imports,
activation requires both $c_p=b_p$ and $b_p=M.p$. The second obligation anchors
the selected definition even when both unqualified contexts resolve $p$ to the
same replacement. Dependency visibility therefore does not establish identity;
the association proof does.
Its inventory includes right-hand-side predicate applications and the empty
heap. This prevents a later import from certifying a different constructor or
callee under the same name. Association also checks the selected memory-rule
module and the elaborated interpretations of the formula's notation: empty
heap, pure embedding, conjunction, separating conjunction, typed existential
quantifiers and each used scalar-storage form. For example, both contexts
capture the function $\lambda(a,v).\;a\mathbin{\#}\mathrm{Int}\mapsto v$
before importing the equation bindings and prove equality of those functions.
This detects a changed storage notation even when the named predicate itself
is unchanged. Certificates lacking this interpretation coverage cannot
authorize execution. Computational field rules require the selected
field notation and layout modules to match the generated target interpretation.
Constants lowered to executor integers have the same integer meaning on the
equation's right-hand side; the original definition must prove that equation.
The generator also accepts ordinary Assertion definitions. A definition
without a constructor match produces one unconditional unfolding clause;
a matching definition uses the same branch representation as a Fixpoint.
Explicit existential binders retain their types and order. Scalar, list and
registered nominal witnesses become fresh logical variables of the corresponding
executor type, while the Rocq equation retains the original typed quantifiers.
Unfolding exposes resources without selecting witness values. Constructor
selection proposals are generated only for matching clauses. Unsupported
types and terms remain explicit normalization barriers.
Nested list patterns are translated as typed constructor trees. Simple local
type aliases are resolved with cycle and ambiguity checks; each nested
constructor must have a unique declared signature with the expected result
type and arity. Binder leaves retain their inferred argument types, while
constructor and nominal-type identities enter the association inventory.
The same tree supplies the strategy pattern, its Rocq equation and any
goal-side model proposal. Implicit existential binders are admitted when an
address or scalar-storage occurrence determines a scalar witness, or a unique
explicitly typed definition determines the type of a direct argument; their
explicitly typed equations still require independent proof. Known aggregate
values cannot be lowered to scalar storage.
Pure logical applications remain opaque unless an admitted computation rule
reduces them. Definitions inside submodules retain their qualified identities
in both the executor strategies and the proof equations.

Witness existence is part of the assertion's semantics. For a witness absent
from the body, the valid equivalence is
\[
  (\exists w:A.\,P(h))\ \Longleftrightarrow\
  (\mathit{inhabited}(A)\land P(h)).
\]
The executor may erase an unused integer or list witness because zero and the
empty list supply inhabitants. A registered nominal type retains its witness
and nonemptiness obligation. Before frontend registration, typed binders are
made syntactically live with the equivalent form
$\exists w:A.\,(w=w\land P)$; this also preserves empty-type quantifiers.
An unresolved inhabitance obligation remains a VC for checked proof.
Nested structure and union fields use the exported target layout. For a root
type $T_0$, symbolic base $a_0=p$, and field sequence $f_1,\ldots,f_k$, the
generator requires an admitted lookup
$(\delta_i,T_i)=\mathit{lookup}(T_{i-1},f_i)$ and constructs
$a_i=a_{i-1}+\delta_i$. Each intermediate type must support the next field
selection; a pointer introduces no implicit dereference. The emitted equation
retains this addition order and checks the original field assertion against
the corresponding scalar resource. The layout bytes and digest are frozen,
and association checks them against the actual execution target. The shared
executor normalizer uses the same admitted field metadata to normalize memory
accesses and owned cells. Integer storage retains its width and signedness.
The existing mechanized path-composition and interval lemmas describe this
address arithmetic; the generator and its connection to frontend metadata
remain part of the implementation trust boundary.
The controller pins source, header, strategy and logical dependency inputs,
executes the selected instance, and compares the resulting layout, import
context and dependency inventory with the preflight. Equation and association
proofs retain separate scopes from the execution VCs. This interface currently
handles one selected instance per invocation; batch activation and certificate
caching remain subsequent integration work.
Normalization must preserve the joint interpretation of $P$ and every term
referenced by the continuation, including return values and call frames.
Assertion equivalence alone is insufficient: eliminating $r$ from
$\exists r.\ r=2 \land \mathit{cell}(p,1)$ while an external return term still
refers to $r$ loses the return constraint. The implementation protects variables
existing before predicate expansion and only eliminates suitable fresh
definition-local witnesses. C arithmetic still emits the ordinary safety
obligations for overflow, division, shifts and memory operations.

\paragraph{Recovering logical models at a checking boundary.}
Expanding a heap predicate can expose the fields needed for execution while
removing the syntactic predicate application used to infer a callee's logical
arguments. Our experimental rule generator reconstructs such arguments by
proposing constructor witnesses on the goal side. For a predicate $H$, a
constructor $C$, and fresh constructor arguments $\bar u$, the refinement is
\[
\frac{P \vdash \exists\bar u.\ t=C(\bar u) \land H(p,C(\bar u))}
     {P \vdash H(p,t)}.
\]
The constructor equality and every resource required by its defining clause
remain obligations. Pure guards guide witness selection; they do not become
new assumptions. Literal guard matching and matching an existing symbolic
equality are both needed because simplification can remove a tautological
guard. Generated Rocq lemmas check constructor substitution and branch
reconstruction. This prototype operates on the boundary's checking copy;
the execution assertion continues independently. Search may remain incomplete,
and the strategy-syntax translation remains trusted.

\paragraph{Proof-producing evaluation of ground integer terms.}
An experimental evaluation service uses Rocq's evaluator on closed logical
terms in their imported definition environment. A computed integer $n$ for
an expression $e$ is a candidate: the service separately generates and checks
the equality $e=n$. An opaque term or a failed check produces no equality.
These equalities can rewrite residual VCs while preserving the original
definitions. The prototype accepts integer-valued applications and arithmetic;
it is not yet connected to every expression encountered during execution.
Evaluation and proof compilation record wall time, CPU time, and peak process
memory separately from the execution run.

\subsection{State-preserving checking boundaries}

For an annotation $A$, boundary processing has the form
\[
  (P,\Omega) \longmapsto
  \bigl(P,\Omega \cup \{\mathit{copy}(P)\vdash\mathit{resolve}(A)\}\bigr).
\]
The goal uses the current program-variable environment and the call's logical
entry bindings. The original $P$ continues unchanged. This rule applies to
supported assertions, visited invariant locations and contract boundaries.
Checking a visited invariant establishes its property at those instance states;
it does not by itself establish an inductive invariant for arbitrary inputs.

\paragraph{Named observation alternatives.}
Proof-only branch-routing commands do not change the executing state. If a
checking boundary has alternatives $A_1,\ldots,A_k$ but its source has no
routing label, its target is their complete disjunction:
$P\vdash\bigvee_{i=1}^{k}\mathit{resolve}(A_i)$, with the observation frame
computed for each alternative. Selecting only the first alternative is not
this query. Existing explicit routing labels retain QCP's selected-case
behavior. The fallback applies on private checking copies; neither a selected
case nor a disjunct becomes an assumption of the continuation.

\paragraph{Local storage in auxiliary observations.}
C variable storage is represented by owned cells in the symbolic assertion.
An auxiliary heap assertion need not repeat every live scalar stack slot.
For a goal $A$, the observation interface identifies unmentioned local scalar
cells through the current program-id/address map and copies their resources
into a frame $F_{P,A}$ on the checking goal. The emitted obligation is
$P\vdash\mathit{resolve}(A)*F_{P,A}$; the continuation remains $P$.
A cell whose address is already described in the goal is not framed again.
This interface retains actual values and types, admits no arbitrary user-heap
frame, and leaves aggregate or opaque local storage as an explicit coverage
barrier. Instantiated call-precondition queries also use the shared expression
and address normalizer on private query copies before resource matching.

Goal-local existential values are resolved before frame selection. Suppose the
goal contains $\exists z.\,\mathit{cell}_t(a,z)*B(z)$ and the current assertion
owns $\mathit{cell}_t(a,v)$ at a mapped local address. The checker chooses $v$
as the goal witness, yielding $\mathit{cell}_t(a,v)*B(v)$ on the query copy.
It requires matching addresses and types, a bare existential variable and no
self-reference in $v$. Every substitution removes one goal existential, so
this witness-selection pass terminates. It then normalizes the goal and
selects its frame. For example, the goal generated by
\texttt{data\_at(q,int,1)} reads a local pointer slot through an existential
value. If the current slot holds $\&i$, choosing that value exposes the
already-described $i$ cell before framing. Each execution branch uses its own
current state; neither entry values nor another branch's witnesses are assumed.

The mechanized observation model gives this interface an exact interpretation.
Let $m$ identify the admitted frame addresses and decompose a heap $h$ into
$h_m$ and $h_{\neg m}$. For an extensional assertion $A$,
\[
 (A*\mathit{exact}(h_m))(h)\ \Longleftrightarrow\ A(h_{\neg m}).
\]
The proof establishes disjoint decomposition and uniqueness of the visible
component, and shows that both stack and visible cells are retained. An extra
visible cell prevents an empty-heap observation from succeeding. Thus local
framing does not make separation logic affine. The abstract mask and typed
cell interpretation remain parameters; linking them to C frontend metadata
is part of the implementation-refinement obligation.
For existential observations, the address mask depends on the witness:
\[
 \exists w.\,(A_w*\mathit{exact}(h_{m_w}))(h)
 \ \Longleftrightarrow\ \exists w.\,A_w(h_{\neg m_w}).
\]
The mechanization proves this equivalence and the sufficiency of a checked
chosen witness. Failure with that witness leaves an obligation; it does not
refute all possible witnesses.

The compatibility flags \texttt{expand} and \texttt{no-expand} both execute
available bodies and loops. The former ignores intermediate annotations by
default, while the latter checks them on state copies. Both always check entry
satisfiability, safety and the target's final Ensure. The execution semantics
therefore remain the same across the two policies.

\subsection{Function calls and loops}

Each call first evaluates its actual arguments using QCP, retaining each
resulting path and its value tuple. The executor allocates fresh parameter and
local storage for the invocation and saves the caller's environment. Passing a
pointer copies its value into a fresh parameter slot; it preserves the identity
of the pointed-to storage and hence aliasing and shared writes. Structure-value
arguments use QCP's typed record-copy rules. The callee's real body is executed
with its own control-flow state, including recursive invocations.

When boundary checking is enabled, QCP's usual contract matching runs on a
state copy to produce Require and frame obligations and bind the callee's
logical inputs. Those entry bindings also determine its return checks. Each
actual return preserves its value, shared heap updates and feasible path.
After typed local cleanup and caller-environment restoration, the executor
checks the instantiated Ensure with its frame and actual return-value relation,
then resumes the caller from the actual returned state. The Ensure never
supplies the caller's execution state. Failed contract binding is reported
explicitly; some current failures still prevent execution and remain coverage
gaps.

We mechanize the frame discipline separately in a scalar-cell reference model.
For already evaluated arguments $v_0,\ldots,v_{m-1}$ and fresh logical identity
$b$, entry installs $H(b+i)=v_i$, maps parameter $i$ to $b+i$, and pushes
$(b,\rho_{\mathrm{caller}})$. Pointer values are copied unchanged. Local
allocation advances the fresh-identity counter. On return, the model retains
the returned value and heap writes, retires callee cells in the allocated
interval, restores $\rho_{\mathrm{caller}}$, and pops the frame. The counter
remains monotone. Thirteen closed Rocq theorems establish argument copying,
alias preservation, shared-write preservation, environment restoration and
local-storage retirement, including nested calls and escaped local addresses.
These are reference semantics for scalar stack cells; correspondence with C
argument evaluation, aggregate layout and dynamic allocation is a separate
implementation obligation.

Loops repeatedly evaluate their conditions and execute their bodies, respecting
\texttt{break}, \texttt{continue} and early return. Each selected invariant
boundary observes the current iteration state. Budget exhaustion records an
incomplete execution and cannot produce a completed-test success. Execution
currently requires available callee bodies; a declaration-only target can
instead perform an explicitly labelled Require-only test.

\subsection{Proof completion and evidence}

The automatic controller produces a run report and immutable VC/dependency
artifacts. QCP's entailment strategies and SMT solver attempt automatic
discharge. Residual obligations are exposed as Rocq statements, with a
separate optional agent workflow for proof construction. A fixed external
checker verifies the submitted proof files against the frozen goals and records
time, memory and available agent token usage.

For the computational field-address candidate, the executor also exports a
target-layout record containing address width, byte order, memory rules,
aggregate identities, and each field's offset, size, alignment and type.
Aggregate references use identities, so recursive pointer types have finite
descriptions. The VC printer and interval checker share the same field-layout
admission test. The layout record travels with the original VC manifest and
is included in the certificate manifest held by the controller. This binds
proof artifacts to the reported interpretation; it does not establish that
legacy imported predicates use that interpretation. Such predicates require
consistent re-elaboration of their address and memory definitions.

The candidate interpretation uses a partial evaluator
$\mathit{addrEval}_{\Lambda}(e)$, where $\Lambda$ resolves type identities,
aggregate sizes and field descriptors. Field selection adds the declared byte
offset; array indexing and pointer arithmetic scale by the resolved element
size. The result contains both the address expression and its type, allowing
nested selections to compose. Unknown fields, ambiguous names and unsupported
address forms have no result. A generated Rocq notation applies this evaluator
while elaborating predicate definitions, retaining symbolic bases and indices
in the resulting arithmetic expression. Thus the predicate's own field
assertions use the same layout as the generated VC. Re-elaboration selects
this explicit interpretation; already compiled predicates retain their
original interpretation. Successful address evaluation establishes neither
ownership nor bounds, which remain separate proof obligations.

Definition lowering uses the same admitted address-term translation for storage
atoms and spatial-predicate arguments. Parentheses delimit application arguments,
so a recursive call carrying a field address retains that term as a single
argument. Known nested paths compose layout offsets; an unknown layout or
implicit pointer traversal leaves the definition unsupported. The controller
checks branch equations separately from program association: a valid equation
over an abstract field interpretation does not authorize a computed address in
execution.

Typed logical terms retain their function identities through lowering. Standard
list append and reverse use polymorphic list interfaces; their generated
certificates also require equality with the standard functions and list
constructors. Program association captures those names before importing the
equation context, rejecting same-named replacements. After ordinary evaluation,
the shared simplifier decides equality of finite integer lists whose type,
constructor shape and elements are known. Unknown elements or tails and exhausted
comparison budgets retain the formula. Five closed theorems establish the
denotation-level comparison contract and preservation of owned heap observations;
closure checks and integer decoding remain implementation-refinement obligations.

Run-level proof completion freezes the automatic report and links each proof
bundle to its original case and VC manifest. The external controller replays
the complete original VC interface, including initial satisfiability, before
updating that case's verdict. Completed executions with accepted certificates
can move from unresolved to accepted; execution limits and errors retain their
original verdicts. A Require-only target retains its specification-only scope,
and ignored annotations remain unchecked. Runs without eligible proof tasks
skip proof construction and replay.

Evidence is tracked by scope. Automatically accepted execution results rely on
the executor and solver trusted base; kernel-checked residual proofs certify
their stated obligations. A whole-instance certificate additionally needs
complete VC coverage and a transitive assumption audit. Existing automatic
proof stubs cannot be counted as checked proofs. This distinction is part of
the experimental data model and proof-workflow evaluation.

The assumption audit is operationally necessary. In the two QCP logical
runtime versions examined, an inverse-field-address axiom equates expression
syntax trees whenever their field addresses coincide. Combined with existing
evaluation and congruence axioms, it permits a proof of false by equating two
distinct syntax constructors. Our adversarial regression compiles this
candidate and then rejects it through the fixed transitive-assumption policy
in both runtimes. Thus, compilation success alone cannot serve as the test
oracle; imported domain axioms must satisfy the certificate admission policy.

\subsection{VC states and test verdicts}

Figure~\ref{fig:result-flow} connects the precondition, intermediate
annotation, and postcondition of one program to their generated obligations.
The precondition asks whether the input domain is inhabited; the remaining
VCs ask whether the selected property holds on the corresponding path state.
Each obligation first goes to ATP, including SMT. A decided result is recorded
immediately. An unresolved obligation goes to the agent and ITP: the agent
assesses the proposition, chooses a proof direction, and submits a proof of
$\phi_i$ or $\neg\phi_i$. ITP checks that proof. If the first direction cannot
be established, the other direction remains available; if neither direction
is established, the result is UNKNOWN.

\begin{figure*}[t]
  \centering
  \resizebox{\textwidth}{!}{\InputIfFileExists{figures/figure-style.tex}{}{\input{figure-style}}
\begin{tikzpicture}[x=1cm,y=1cm,font=\footnotesize,text=tsInk]
\path[tsPanel,fill=tsProof!3] (0,0) rectangle (8.05,-6.6);
\path[tsPanel,fill=tsExec!3] (8.55,0) rectangle (17.35,-6.6);
\node[tsHead] at (.18,-.35) {(a) Evidence for one obligation};
\node[tsHead] at (8.73,-.35) {(b) Evidence for one instance};
\node[align=center] at (4.02,-.97)
 {$\phi_j\equiv\forall z.\,P_j(z)\Rightarrow A_j(z)$\\
 SMT query:\ $\exists z.\,P_j(z)\land\neg A_j(z)$};
\node[anchor=west] (unsat) at (.25,-1.8) {\textbf{UNSAT}};
\node[anchor=west] (sat) at (.25,-2.65) {\textbf{SAT}};
\node[anchor=west] (unknown) at (.25,-3.5) {\textbf{UNKNOWN}};
\node[tsCard,text width=2.35cm,minimum height=.45cm] (model) at (3.4,-2.65) {validate\\original model};
\node[tsGood,text width=1.4cm] (valid) at (6.85,-1.8) {VALID};
\node[tsBad,text width=1.4cm] (invalid) at (6.85,-2.65) {INVALID};
\node[tsOpen,text width=1.4cm] (open) at (6.85,-3.5) {unresolved};
\draw[tsEvidence] (unsat.east)--node[tsLabel,above]{sound encoding}(valid.west);
\draw[tsEvidence] (sat.east)--(model.west);
\draw[tsEvidence] (model.east)--node[tsLabel,above]{checked}(invalid.west);
\draw[tsEvidence] (unknown.east)--node[tsLabel,below]{or unsupported}(open.west);
\draw[tsEvidence] (model.south)--(3.4,-3.12)-|node[tsLabel,near start,above]{unconfirmed}(open.north);
\node[tsCard,draw=tsProof,text width=6.8cm,minimum height=.58cm] (itp) at (4.02,-4.65)
 {\textbf{ITP}\quad construct candidate $p$\ $\longrightarrow$\ kernel + audit};
\draw[tsEvidence] (open.south)--(6.85,-4.1)-|(itp.north);
\node[align=center,text width=2.2cm] (pos) at (1.4,-5.65)
 {$p:\phi_j$\\\textcolor{tsValid}{VALID}};
\node[align=center,text width=2.2cm] (neg) at (4.02,-5.65)
 {$p:\neg\phi_j$\\\textcolor{tsInvalid}{INVALID}};
\node[align=center,text width=2.2cm] (none) at (6.65,-5.65)
 {no checked proof\\\textcolor{tsUnknown}{UNKNOWN}};
\foreach \n in {pos,neg,none}{\draw[tsEvidence] (itp.south -| \n.north)--(\n.north);}
\node[font=\scriptsize,text=tsUnknown] at (4.02,-6.3) {A failed proof attempt supplies no refutation.};
\node[align=center,font=\footnotesize] at (12.95,-1.0)
 {\textbf{Schematic VC ledger}\quad $\checkmark$ valid\quad $\times$ invalid\quad ? unresolved};
\node[inner sep=0pt] at (12.95,-2.45) {\renewcommand{\arraystretch}{1.3}
 \begin{tabular}{@{}c|ccc|l@{}}
 path & safety & assertion & return & path verdict\\\hline
 $\pi_1$ & $\checkmark$ & $\checkmark$ & $\checkmark$ & \textcolor{tsValid}{PASS}$^{a}$\\
 $\pi_2$ & $\checkmark$ & $\times$ & ? & \textcolor{tsInvalid}{FAIL}$^{b}$\\
 $\pi_3$ & $\checkmark$ & ? & $\checkmark$ & \textcolor{tsUnknown}{UNKNOWN}
 \end{tabular}};
\node[align=left,anchor=west,text width=7.8cm] at (8.95,-3.65)
 {$^{a}$ Feasible path; execution and selected checks complete.\\
 $^{b}$ The refuted check has a reachable violating state.};
\draw[tsEvidence] (12.95,-4.05)--(12.95,-4.4);
\node[tsCard,text width=7.5cm,minimum height=1.05cm,align=left] at (12.95,-5.05)
 {\textbf{Instance verdict}\quad (first applicable rule)\\[3pt]
 empty entry or some path FAIL $\Rightarrow$ \textbf{FAIL}\\
 nonempty entry + coverage + all paths PASS $\Rightarrow$ \textbf{PASS}\\
 otherwise $\Rightarrow$ \textbf{UNKNOWN}};
\node[align=center,text width=7.85cm,font=\scriptsize] at (12.95,-6.15)
 {Proven-infeasible paths are discarded. Empty entry is an invalid test,\\
 reported separately from a reachable execution violation.};
\draw[tsEvidence] (8.08,-2.45)--(8.52,-2.45);
\end{tikzpicture}
}
  \caption{From program positions to a test result. The precondition generates
  a satisfiability query; annotations and the postcondition generate validity
  queries. Every query goes first to ATP and then, if unresolved, to Agent+ITP.
  Checked positive and negative proofs determine the corresponding true and
  false results; unsuccessful attempts in both directions leave UNKNOWN.
  Three-valued conjunction aggregates the individual results.}
  \label{fig:result-flow}
\end{figure*}

Each selected obligation has state \valid{}, \invalid{}, or \unknown{}.
The SMT interface can check $P\land\neg A$ for $P\vdash A$ and provide a
counterexample candidate. Path infeasibility is a separate satisfiability
query on $P$; its UNSAT result prunes the path. A candidate only establishes
\invalid{} after the negative obligation and reachable path witness are
admitted. This separation prevents an inconsistent antecedent from being
reported as an invalid VC.

The implementation checks violations against a private copy of the original
path assertion, before entailment strategies transform it. Its admitted
satisfiability procedure retains pure constraints and heap ownership. It can
establish a failed pure condition or a conflicting value in an owned scalar
cell. Disjunctive observations are checked jointly. Conditions depending on
unresolved goal-local witnesses cannot justify rejection of all witnesses;
unsupported formulas remain \unknown{}. A positive ITP proof changes a VC to
\valid{}. For residual VCs, the agent may try both $\mathrm{VC}$ and
$\neg\mathrm{VC}$, and ITP checks either submitted proof. Failed attempts in
both directions leave the VC unresolved.

Entry results are SAT, UNSAT, or UNKNOWN. To aggregate all obligations,
SAT and VALID map to $\mathsf T$, UNSAT and INVALID to $\mathsf F$, and UNKNOWN
to $\mathsf U$. Three-valued conjunction yields FAIL if any component is
$\mathsf F$, PASS if every component is $\mathsf T$, and UNKNOWN otherwise.
An incomplete execution contributes $\mathsf U$. An UNSAT entry stops execution
and reports FAIL with an unsatisfiable-entry reason. The executor continues
from the original state after an observation, so a failing annotation does
not replace the state used by later checks.

\paragraph{Regression acceptance.}
End-to-end regression checks compare the final test verdict with the case's
expected behavior. Positive cases require PASS. For negative cases, the weak
criterion accepts FAIL or UNKNOWN; the strong criterion requires FAIL.
A negative case returning PASS fails both criteria. Reports record the chosen
criterion and retain the original execution verdict, so an UNKNOWN result
accepted under the weak criterion remains distinguishable from a confirmed
negative result.

Certificate replay records proof polarity separately from VC state. A checked
negative proof of the original entry-existence obligation establishes UNSAT.
A checked negative proof of a selected path VC establishes \invalid{} when
the path is admitted as reachable. Archived negative proofs without that
reachability evidence remain proof evidence rather than an execution \fail{}.

\paragraph{Root instances and reusable contracts.}
The substitution $\theta$ specializes the root precondition, not the reusable
contract of every invocation of the same function. At a recursive call, actual
arguments and the current heap determine a fresh contract instance. Reusing
root bindings there would, for example, incorrectly require a decreasing
counter to retain its initial value. The frontend preserves the original
contracts and emits a fresh named root specification selected by QCP's entry
option. Its position preserves QCP's default call-contract selection, and
explicitly named calls continue to reference the original names. C bodies and
call expressions are unchanged. Declaration-only checks retain a single
initial-domain obligation.


\paragraph{External definition closure.}
Certificate preparation stages locally resolved logical dependencies in
import order, including \texttt{Require Import}, \texttt{Require Export} and
plain \texttt{Require}. Nested comments and string literals are masked before
collecting imports. Source and dependency hashes are frozen before proof
production; successful compilation and independent replay remain necessary
for accepting the resulting closure.


\paragraph{Explicit library interpretation.}
A preparation utility accepts target-layout metadata and selected logical
source modules. It preserves their names, declarations, and proofs, inserting
the computational address notation at a unique existing scope boundary.
Comments and strings are excluded when locating that boundary. The bundle
records original and transformed source hashes and starts in an unchecked
state. Compilation and the normal certificate controller provide the subsequent
acceptance gates. Selecting modules for this interpretation is explicit;
preparation does not certify equivalence to previously compiled abstract
address definitions or infer a reinterpretation of all imported dependencies.
